% ch3-linear-systems.tex — Chapter 3: Systems of Linear Equations and Gaussian Elimination
% Systems of linear equations, row reduction, echelon forms, rank,
% Rouché–Capelli theorem, structure of solutions, applications.

\chapter{Systems of Linear Equations and Gaussian Elimination}%
\label{ch:linear-systems}

% ============================================================
% Motivating introduction
% ============================================================

Systems of linear equations are among the oldest and most ubiquitous
objects in mathematics. They appear whenever several quantities are
linked by proportional relationships: currents in an electrical
network must satisfy Kirchhoff's laws, prices in a multi-market
economic equilibrium are determined by supply-and-demand balances,
forces on a bridge truss must cancel for the structure to remain
in static equilibrium. Even the modern theory of machine learning
ultimately reduces the training of linear models to the solution of
(often very large) linear systems.

The algorithmic heart of the subject is \emph{Gaussian elimination}%
\index{Gaussian elimination}, a systematic procedure that transforms
any system into an equivalent one whose solutions can be read off by
back-substitution. Despite its simplicity, the algorithm is
remarkably powerful: it simultaneously tells us whether a system has
solutions, how many free parameters the solution set carries, and
provides an explicit parametric description.

In this chapter we formalise systems of linear equations, introduce
the matrix language~$Ax=b$, develop row reduction to (reduced) echelon
form, define the rank of a matrix, and prove the
Rouch\'{e}--Capelli theorem that characterises consistency. We then
establish the fundamental structure theorem: the general solution of a
consistent system is a translate of the solution space of the
associated homogeneous system. Throughout, we provide geometric
interpretations and worked examples.

\medskip
Throughout this chapter, $\K$ denotes a field (typically $\R$ or $\C$),
and we freely use results from 
ef{ch:vector-spaces} and

ef{ch:linear-maps}.


% ============================================================
\section{Systems of linear equations}\label{sec:linear-systems}
% ============================================================

\begin{definition}{System of linear equations}{def:linear-system}
\index{system of linear equations}\index{linear system}
A \emph{system of $m$ linear equations in $n$ unknowns} over a
field~$\K$ is a collection of equations
\[
  \left\{
  \begin{aligned}
    a_{11}\,x_1 + a_{12}\,x_2 + \cdots + a_{1n}\,x_n &= b_1,\\
    a_{21}\,x_1 + a_{22}\,x_2 + \cdots + a_{2n}\,x_n &= b_2,\\
    &\;\;\vdots \\
    a_{m1}\,x_1 + a_{m2}\,x_2 + \cdots + a_{mn}\,x_n &= b_m,
  \end{aligned}
  \right.
\]
where $a_{ij}\in\K$ are the \emph{coefficients}\index{coefficient}
and $b_i\in\K$ are the \emph{right-hand sides} (or \emph{constant
terms}). A vector $(s_1,\dots,s_n)\in\K^n$ is a \emph{solution} if
substituting $x_j = s_j$ satisfies every equation simultaneously. The
\emph{solution set}\index{solution set} is
\[
  \cS \deff \setcond{x\in\K^n}{Ax = b}.
\]
\end{definition}

\subsection{Matrix form}\label{subsec:matrix-form}
\index{matrix form of a linear system}

The system above can be written compactly as a single matrix equation
\begin{equation}\label{eq:Axb}
  Ax = b,
\end{equation}
where
\[
  A = \begin{pmatrix}
    a_{11} & a_{12} & \cdots & a_{1n}\\
    a_{21} & a_{22} & \cdots & a_{2n}\\
    \vdots & \vdots & \ddots & \vdots\\
    a_{m1} & a_{m2} & \cdots & a_{mn}
  \end{pmatrix}
  \in \Mat_{m,n}(\K),
  \qquad
  x = \begin{pmatrix} x_1\\ x_2\\ \vdots\\ x_n \end{pmatrix}
  \in \K^n,
  \qquad
  b = \begin{pmatrix} b_1\\ b_2\\ \vdots\\ b_m \end{pmatrix}
  \in \K^m.
\]
The matrix $A$ is called the \emph{coefficient matrix}%
\index{coefficient matrix}. The \emph{augmented matrix}%
\index{augmented matrix} is the $m\times(n+1)$ matrix
\begin{equation}\label{eq:augmented}
  (A \mid b) \deff
  \left(\begin{array}{cccc|c}
    a_{11} & a_{12} & \cdots & a_{1n} & b_1\\
    a_{21} & a_{22} & \cdots & a_{2n} & b_2\\
    \vdots & \vdots & \ddots & \vdots & \vdots\\
    a_{m1} & a_{m2} & \cdots & a_{mn} & b_m
  \end{array}\right).
\end{equation}

\begin{remark}{Column picture}{rem:column-picture}
\index{column picture}
The equation $Ax = b$ can also be read as asking whether $b$ lies in
the column space of~$A$: writing $A = (\vec{a}_1 \mid \vec{a}_2 \mid
\cdots \mid \vec{a}_n)$, the system is consistent if and only if
\[
  b = x_1\,\vec{a}_1 + x_2\,\vec{a}_2 + \cdots + x_n\,\vec{a}_n
\]
for some scalars $x_j\in\K$; that is, $b\in\im(A)$ where we view $A$
as a linear map $\K^n\to\K^m$.
\end{remark}


% ============================================================
\section{Homogeneous and non-homogeneous systems}%
\label{sec:homogeneous}
% ============================================================

\begin{definition}{Homogeneous system}{def:homogeneous}
\index{homogeneous system}\index{non-homogeneous system}
The system $Ax = b$ is called \emph{homogeneous} if $b = \vzero$, and
\emph{non-homogeneous} (or \emph{inhomogeneous}) otherwise.
Given a (possibly non-homogeneous) system $Ax = b$, the system
$Ax = \vzero$ is called its \emph{associated homogeneous system}.
\end{definition}

\begin{proposition}{The solution set of a homogeneous system is a subspace}%
{prop:homog-subspace}
\index{kernel!of a matrix}
Let $A\in\Mat_{m,n}(\K)$. The solution set of $Ax=\vzero$ is a
vector subspace of~$\K^n$. It coincides with $\Ker(A)$, the kernel
of the linear map $x\mapsto Ax$.
\end{proposition}

\begin{proof}
The zero vector satisfies $A\vzero = \vzero$, so $\cS_0\neq\emptyset$.
If $u,v\in\cS_0$ and $\lambda\in\K$, then
$A(\lambda u + v) = \lambda Au + Av = \lambda\vzero + \vzero = \vzero$,
so $\lambda u + v\in\cS_0$.
By the subspace criterion (
ef{ch:vector-spaces}), $\cS_0$ is a
subspace of~$\K^n$. By definition,
$\cS_0 = \setcond{x\in\K^n}{Ax = \vzero} = \Ker(A)$.
\end{proof}

\begin{proposition}{Affine structure of the solution set}%
{prop:affine-structure}
\index{affine subspace}
Let $Ax = b$ be a consistent system (i.e.\ $\cS\neq\emptyset$), and
let $x_0\in\cS$ be any particular solution. Then
\begin{equation}\label{eq:affine-sol}
  \cS = x_0 + \Ker(A) \deff \setcond{x_0 + h}{h \in \Ker(A)}.
\end{equation}
In particular, $\cS$ is an \emph{affine subspace}\index{affine subspace}
of~$\K^n$ of dimension $\dim\Ker(A)$.
\end{proposition}

\begin{proof}
($\supseteq$) If $h\in\Ker(A)$, then $A(x_0+h) = Ax_0 + Ah = b +
\vzero = b$, so $x_0+h\in\cS$.

($\subseteq$) If $x\in\cS$, set $h\deff x - x_0$. Then
$Ah = Ax - Ax_0 = b - b = \vzero$, so $h\in\Ker(A)$ and
$x = x_0 + h$.
\end{proof}


% ============================================================
\section{Elementary row operations}\label{sec:row-operations}
% ============================================================

\begin{definition}{Elementary row operations}{def:ero}
\index{elementary row operations}\index{row operations}
The following three operations on the rows of a matrix are called
\emph{elementary row operations} (EROs):
\begin{enumerate}[label=(E\arabic*)]
  \item\label{ero:swap}
    \textbf{Row interchange:} Swap rows $i$ and~$j$
    \quad ($R_i \leftrightarrow R_j$).
  \item\label{ero:scale}
    \textbf{Row scaling:} Multiply row $i$ by a non-zero scalar
    $\lambda\in\K^*$ \quad ($R_i \leftarrow \lambda\,R_i$).
  \item\label{ero:add}
    \textbf{Row replacement:} Add $\lambda$ times row $j$ to row $i$
    \quad ($R_i \leftarrow R_i + \lambda\,R_j$), where $i\neq j$.
\end{enumerate}
\end{definition}

\begin{proposition}{EROs preserve the solution set}{prop:ero-preserve}
If the augmented matrix $(A'\mid b')$ is obtained from $(A\mid b)$ by
an elementary row operation, then the systems $Ax=b$ and $A'x=b'$
have the same solution set.
\end{proposition}

\begin{proof}
Each ERO is reversible: \ref{ero:swap} is its own inverse;
\ref{ero:scale} is reversed by scaling by $\lambda^{-1}$;
\ref{ero:add} is reversed by subtracting $\lambda$ times row~$j$.
Hence any solution of one system is a solution of the other.
More precisely, each ERO amounts to left-multiplication by an
invertible \emph{elementary matrix}\index{elementary matrix}
$E\in\GL_m(\K)$. Since $Ax = b$ if and only if $E(Ax) = Eb$, i.e.\
$(EA)x = Eb$, the solution set is unchanged.
\end{proof}

\begin{definition}{Row equivalence}{def:row-equiv}
\index{row equivalence}
Two matrices are \emph{row equivalent} if one can be obtained from
the other by a finite sequence of elementary row operations. We write
$A \sim A'$.
\end{definition}


% ============================================================
\section{Echelon forms}\label{sec:echelon}
% ============================================================

\begin{definition}{Row echelon form}{def:ref}
\index{row echelon form}\index{echelon form}
A matrix is in \emph{row echelon form} (REF) if:
\begin{enumerate}[label=(\roman*)]
  \item All zero rows are at the bottom.
  \item The first non-zero entry in each non-zero row (called the
    \emph{pivot}\index{pivot}) is strictly to the right of the pivot
    in the row above.
\end{enumerate}
\end{definition}

\begin{definition}{Reduced row echelon form}{def:rref}
\index{reduced row echelon form}\index{RREF}
A matrix is in \emph{reduced row echelon form} (RREF) if it is in REF
and, additionally:
\begin{enumerate}[label=(\roman*),resume]
  \item Every pivot is equal to~$1$.
  \item Each pivot is the only non-zero entry in its column.
\end{enumerate}
\end{definition}

\begin{example}{Echelon forms}{ex:echelon-shapes}
The following matrix is in row echelon form (pivots are boxed):
\[
  \begin{pmatrix}
    \boxed{2} & 3 & -1 & 7\\
    0 & \boxed{5} & 4 & 2\\
    0 & 0 & 0 & \boxed{1}\\
    0 & 0 & 0 & 0
  \end{pmatrix}.
\]
Its reduced row echelon form is
\[
  \begin{pmatrix}
    \boxed{1} & 0 & * & 0\\
    0 & \boxed{1} & * & 0\\
    0 & 0 & 0 & \boxed{1}\\
    0 & 0 & 0 & 0
  \end{pmatrix},
\]
where the $*$ entries are determined by the original entries.
\end{example}

\begin{theorem}{Existence and uniqueness of RREF}{thm:rref-unique}
\index{RREF!uniqueness}
Every matrix $A\in\Mat_{m,n}(\K)$ is row equivalent to a unique
matrix in reduced row echelon form.
\end{theorem}

\begin{proof}
\emph{Existence} follows from the Gaussian elimination algorithm
described in 
ef{sec:gauss-elim}, which produces a REF, followed
by Gauss--Jordan elimination (
ef{sec:gauss-jordan}), which
produces the RREF.

\emph{Uniqueness.}
Suppose $R$ and $R'$ are two RREFs row equivalent to~$A$. Then
$R$ and $R'$ are row equivalent to each other, so they define the
same system of homogeneous equations $Rx=\vzero$ and $R'x=\vzero$
(by 
ef{prop:prop:ero-preserve}). In particular, their solution sets
coincide: $\Ker(R) = \Ker(R')$.

We show $R = R'$ column by column. Let $j$ be any column index.
If $j$ is a pivot column of~$R$, say with pivot in row~$i$, then
the $i$-th equation of $Rx = \vzero$ reads $x_j = -\sum_{k}
r_{ik}\,x_k$ where the sum is over the non-pivot (free) variables.
Since $\Ker(R) = \Ker(R')$, column~$j$ must also be a pivot column
of~$R'$ in the same row, and the coefficients must agree. Similarly
for non-pivot columns. Hence $R = R'$.
\end{proof}


% ============================================================
\section{Gaussian elimination}\label{sec:gauss-elim}
% ============================================================

\index{Gaussian elimination!algorithm}
Gaussian elimination transforms a matrix into row echelon form using
elementary row operations. The algorithm proceeds as follows.

\medskip
\noindent\textbf{Algorithm} (Gaussian elimination).
\begin{enumerate}[label=\textbf{Step \arabic*.}]
  \item Start with column $j=1$ and row $i=1$.
  \item \emph{Pivot selection.} In column $j$, find a non-zero entry
    in rows $i, i+1, \dots, m$. If none exists, increment $j$ and
    repeat. If $j > n$, stop.
  \item \emph{Swap.} If the non-zero entry is in row $k\neq i$,
    swap $R_i \leftrightarrow R_k$.
  \item \emph{Elimination.} For each row $\ell = i+1, \dots, m$,
    perform $R_\ell \leftarrow R_\ell -
    \frac{a_{\ell j}}{a_{ij}}\,R_i$ to create zeros below the pivot.
  \item Increment $i$ and $j$; go to Step~2.
\end{enumerate}

\begin{example}{Gaussian elimination: a $3\times 4$ system}%
{ex:gauss-3x4}
Solve the system
\[
  \left\{
  \begin{aligned}
    x_1 + 2x_2 - x_3 + 3x_4 &= 5,\\
    2x_1 + 5x_2 + x_3 + 2x_4 &= 8,\\
    x_1 + 3x_2 + 2x_3 - x_4 &= 1.
  \end{aligned}
  \right.
\]
The augmented matrix and its row reduction:
\begin{align*}
  &\left(\begin{array}{cccc|c}
    1 & 2 & -1 & 3 & 5\\
    2 & 5 & 1 & 2 & 8\\
    1 & 3 & 2 & -1 & 1
  \end{array}\right)
  \xrightarrow[\substack{R_3 \leftarrow R_3 - R_1}]{R_2 \leftarrow R_2 - 2R_1}
  \left(\begin{array}{cccc|c}
    1 & 2 & -1 & 3 & 5\\
    0 & 1 & 3 & -4 & -2\\
    0 & 1 & 3 & -4 & -4
  \end{array}\right)\\[6pt]
  &\xrightarrow{R_3 \leftarrow R_3 - R_2}
  \left(\begin{array}{cccc|c}
    1 & 2 & -1 & 3 & 5\\
    0 & 1 & 3 & -4 & -2\\
    0 & 0 & 0 & 0 & -2
  \end{array}\right).
\end{align*}
The last row reads $0 = -2$, a contradiction. Hence the system is
\textbf{inconsistent}\index{inconsistent system}: $\cS = \emptyset$.
\end{example}

\begin{example}{A consistent system with free variables}%
{ex:gauss-free}
Solve the system
\[
  \left\{
  \begin{aligned}
    x_1 + 2x_2 + x_3 &= 4,\\
    2x_1 + 5x_2 + 4x_3 &= 11,\\
    x_1 + 3x_2 + 3x_3 &= 7.
  \end{aligned}
  \right.
\]
Row reduction of the augmented matrix:
\begin{align*}
  &\left(\begin{array}{ccc|c}
    1 & 2 & 1 & 4\\
    2 & 5 & 4 & 11\\
    1 & 3 & 3 & 7
  \end{array}\right)
  \xrightarrow[\substack{R_3 \leftarrow R_3 - R_1}]{R_2 \leftarrow R_2 - 2R_1}
  \left(\begin{array}{ccc|c}
    1 & 2 & 1 & 4\\
    0 & 1 & 2 & 3\\
    0 & 1 & 2 & 3
  \end{array}\right)
  \xrightarrow{R_3 \leftarrow R_3 - R_2}
  \left(\begin{array}{ccc|c}
    1 & 2 & 1 & 4\\
    0 & 1 & 2 & 3\\
    0 & 0 & 0 & 0
  \end{array}\right).
\end{align*}
The pivot columns are $1$ and $2$; the variable $x_3$ is free
\index{free variable}. Back-substitution gives:
\[
  x_2 = 3 - 2x_3, \qquad x_1 = 4 - 2x_2 - x_3 = -2 + 3x_3.
\]
Setting $x_3 = t \in \K$, the solution set is
\[
  \cS = \set{\begin{pmatrix} -2+3t\\ 3-2t\\ t \end{pmatrix}
  : t\in\K}
  = \begin{pmatrix} -2\\3\\0 \end{pmatrix}
  + t\begin{pmatrix} 3\\-2\\1 \end{pmatrix},
  \quad t\in\K.
\]
This is a line in $\K^3$: a particular solution plus the kernel
$\Ker(A) = \Span\!\left(\begin{pmatrix}3\\-2\\1\end{pmatrix}\right)$.
\end{example}


% ============================================================
\section{Gauss--Jordan elimination}\label{sec:gauss-jordan}
% ============================================================

\index{Gauss--Jordan elimination}
After reaching row echelon form, \emph{Gauss--Jordan elimination}
continues with back-elimination to reach the reduced row echelon form.

\medskip
\noindent\textbf{Algorithm} (Gauss--Jordan: from REF to RREF).
\begin{enumerate}[label=\textbf{Step \arabic*.}]
  \item Starting from the \emph{bottom-most} pivot, scale its row so
    the pivot becomes~$1$.
  \item Use this pivot to eliminate all entries \emph{above} it in the
    same column.
  \item Move to the next pivot above and repeat.
\end{enumerate}

\begin{example}{Gauss--Jordan continuation}{ex:gauss-jordan}
Continuing from 
ef{ex:ex:gauss-free}, we had the REF
\[
  \left(\begin{array}{ccc|c}
    1 & 2 & 1 & 4\\
    0 & 1 & 2 & 3\\
    0 & 0 & 0 & 0
  \end{array}\right).
\]
The pivots are already~$1$. Eliminate above the second pivot:
\[
  \xrightarrow{R_1 \leftarrow R_1 - 2R_2}
  \left(\begin{array}{ccc|c}
    1 & 0 & -3 & -2\\
    0 & 1 & 2 & 3\\
    0 & 0 & 0 & 0
  \end{array}\right).
\]
This is the RREF. The solution can be read off directly:
$x_1 = -2 + 3x_3$, $x_2 = 3 - 2x_3$, confirming the result of

ef{ex:ex:gauss-free}.
\end{example}


% ============================================================
\section{Rank of a matrix}\label{sec:matrix-rank}
% ============================================================

\begin{definition}{Rank via echelon form}{def:rank-echelon}
\index{rank!of a matrix}
The \emph{rank} of a matrix $A\in\Mat_{m,n}(\K)$, denoted
$\rank(A)$, is the number of pivots in any row echelon form of~$A$.
\end{definition}

\begin{proposition}{Rank is well-defined}{prop:rank-well-defined}
The rank does not depend on which row echelon form is chosen: it
equals the number of pivots in the unique RREF of~$A$.
\end{proposition}

\begin{proof}
By 
ef{thm:thm:rref-unique}, the RREF is unique. Since the number of
pivots in any REF equals the number of pivots in the RREF (the
Gauss--Jordan steps do not create or destroy pivots), the rank is
well-defined.
\end{proof}

\begin{proposition}{Rank equals column rank equals row rank}%
{prop:rank-equiv}
\index{rank!column rank}\index{rank!row rank}
For any matrix $A\in\Mat_{m,n}(\K)$, the following are equal:
\begin{enumerate}[label=(\roman*)]
  \item The number of pivots in an echelon form of~$A$.
  \item $\dim(\im A)$, i.e.\ the dimension of the column space
    of~$A$.
  \item The dimension of the row space of~$A$.
\end{enumerate}
In particular, $\rank(A) = \rank(\transpose{A})$.
\end{proposition}

\begin{proof}
(i)$=$\,(ii): The pivot columns of~$A$ form a basis for the column
space $\im(A)$ (the corresponding columns of the
\emph{original}~$A$, selected via the pivot positions of the RREF,
are linearly independent and Span~$\im(A)$). Thus
$\dim(\im A)$ equals the number of pivots.

(i)$=$\,(iii): Row operations do not change the row space (each new
row is a linear combination of the old rows, and the operation is
reversible). The non-zero rows of a REF are linearly independent
(by the staircase structure), so they form a basis for the row space.
Their count equals the number of pivots.
\end{proof}

\begin{remark}{Rank and the rank--nullity theorem}{rem:rank-nullity-systems}
\index{rank--nullity theorem}
The rank--nullity theorem (
ef{ch:linear-maps}) applied to the
linear map $x\mapsto Ax$ gives
\begin{equation}\label{eq:rank-nullity}
  \rank(A) + \dim\Ker(A) = n,
\end{equation}
where $n$ is the number of columns (i.e.\ unknowns). Thus the number
of free variables in the general solution equals $n - \rank(A)$.
\end{remark}


% ============================================================
\section{The Rouch\'{e}--Capelli theorem}\label{sec:rouche-capelli}
% ============================================================

The central consistency criterion is the following classical result.

\begin{theorem}{Rouch\'{e}--Capelli theorem}%
{thm:rouche-capelli}
\index{Rouche-Capelli theorem@Rouch\'{e}--Capelli theorem}
\index{consistency criterion}
The system $Ax = b$ is consistent (i.e.\ has at least one solution)
if and only if
\begin{equation}\label{eq:rouche-capelli}
  \rank(A) = \rank(A\mid b).
\end{equation}
When this condition holds, the dimension of the solution set is
$n - \rank(A)$, where $n$ is the number of unknowns. In particular:
\begin{itemize}
  \item If $\rank(A) = n$, the solution is unique.
  \item If $\rank(A) < n$, there are infinitely many solutions
    (assuming $\K$ is infinite), parametrised by $n - \rank(A)$ free
    variables.
\end{itemize}
\end{theorem}

\begin{proof}
Write $r = \rank(A)$ and $r' = \rank(A\mid b)$. Since every column
of~$A$ is also a column of $(A\mid b)$, we always have $r \leq r'
\leq r+1$.

\emph{($\Rightarrow$)}
Suppose $\cS\neq\emptyset$, say $x_0\in\cS$. Then $b = Ax_0$ is a
linear combination of the columns of~$A$, so $b$ lies in the column
space of~$A$. Adding the column~$b$ to~$A$ does not increase the
dimension of the column space:
$\rank(A\mid b) = \dim\Span(\text{columns of } A, b) =
\dim\Span(\text{columns of } A) = \rank(A)$.
Hence $r' = r$.

\emph{($\Leftarrow$)}
Suppose $\rank(A) = \rank(A\mid b)$. Then the column spaces satisfy
$\im(A\mid b) = \im(A)$, which means $b$ is a linear combination of
the columns of~$A$. That is, there exists $x_0\in\K^n$ with
$Ax_0 = b$, so the system is consistent.

The dimension statement follows from 
ef{prop:prop:affine-structure} and

ef{eq:rank-nullity}: $\dim\cS = \dim\Ker(A) = n - \rank(A)$.
\end{proof}

\begin{corollary}{Homogeneous systems always have non-trivial solutions
when $n > m$}{cor:homog-nontrivial}
\index{homogeneous system!non-trivial solution}
If $A\in\Mat_{m,n}(\K)$ with $n > m$ (more unknowns than equations),
then $Ax = \vzero$ has a non-trivial solution.
\end{corollary}

\begin{proof}
We have $\rank(A) \leq m < n$, so $\dim\Ker(A) = n - \rank(A) \geq 1$.
\end{proof}


% ============================================================
\section{Structure of the solution set}\label{sec:solution-structure}
% ============================================================

We now state and prove the complete structure theorem, synthesising
the results of the preceding sections.

\begin{theorem}{Structure of the general solution}{thm:solution-structure}
\index{general solution}\index{solution set!structure}
Let $A\in\Mat_{m,n}(\K)$ and $b\in\K^m$, and suppose $Ax=b$ is
consistent. Let $x_p\in\K^n$ be any particular solution. Then:
\begin{enumerate}[label=(\roman*)]
  \item $\Ker(A) = \setcond{x\in\K^n}{Ax = \vzero}$ is a vector
    subspace of $\K^n$ of dimension $n - \rank(A)$.
  \item The general solution of $Ax = b$ is
    \begin{equation}\label{eq:gen-sol}
      \cS = \setcond{x_p + h}{h\in\Ker(A)} = x_p + \Ker(A).
    \end{equation}
  \item If $\set{h_1, \dots, h_k}$ is a basis of $\Ker(A)$ (where
    $k = n - \rank(A)$), then every solution can be written uniquely as
    \begin{equation}\label{eq:parametric-sol}
      x = x_p + t_1 h_1 + t_2 h_2 + \cdots + t_k h_k,
      \qquad t_1, \dots, t_k \in \K.
    \end{equation}
\end{enumerate}
\end{theorem}

\begin{proof}
(i) follows from 
ef{prop:prop:homog-subspace} and

ef{eq:rank-nullity}.

(ii) is exactly 
ef{prop:prop:affine-structure}.

(iii) Since $\set{h_1,\dots,h_k}$ is a basis of $\Ker(A)$, every
$h\in\Ker(A)$ can be written uniquely as $h = t_1 h_1 + \cdots +
t_k h_k$. Combined with~(ii), every $x\in\cS$ has the unique
representation $x = x_p + t_1 h_1 + \cdots + t_k h_k$.
The uniqueness of the $t_i$ follows from the linear independence of
the $h_i$: if $x_p + \sum t_i h_i = x_p + \sum t_i' h_i$, then
$\sum (t_i - t_i') h_i = \vzero$, forcing $t_i = t_i'$ for all~$i$.
\end{proof}


% ============================================================
\section{Parametric representation of solutions}%
\label{sec:parametric}
% ============================================================

\index{parametric form}
When solving a system $Ax = b$ by Gaussian elimination, the RREF
identifies which variables are \emph{pivot variables}\index{pivot variable}
(corresponding to pivot columns) and which are \emph{free variables}%
\index{free variable} (corresponding to non-pivot columns). A
parametric representation is obtained by:
\begin{enumerate}
  \item Expressing each pivot variable in terms of the free variables.
  \item Assigning arbitrary parameters $t_1, t_2, \dots$ to the free
    variables.
  \item Writing the general solution as a vector.
\end{enumerate}

\begin{example}{Parametric form with two free variables}%
{ex:parametric-two-free}
Consider the system with RREF
\[
  \left(\begin{array}{ccccc|c}
    1 & 3 & 0 & -2 & 0 & 4\\
    0 & 0 & 1 & 5 & 0 & -1\\
    0 & 0 & 0 & 0 & 1 & 2\\
    0 & 0 & 0 & 0 & 0 & 0
  \end{array}\right).
\]
Pivot variables: $x_1, x_3, x_5$.\quad Free variables: $x_2 = s$,
$x_4 = t$. Reading off:
\[
  x_1 = 4 - 3s + 2t, \quad x_3 = -1 - 5t, \quad x_5 = 2.
\]
The parametric form is
\[
  \begin{pmatrix} x_1\\ x_2\\ x_3\\ x_4\\ x_5 \end{pmatrix}
  =
  \underbrace{\begin{pmatrix} 4\\ 0\\ -1\\ 0\\ 2 \end{pmatrix}}_{x_p}
  + s\underbrace{\begin{pmatrix} -3\\ 1\\ 0\\ 0\\ 0 \end{pmatrix}}_{h_1}
  + t\underbrace{\begin{pmatrix} 2\\ 0\\ -5\\ 1\\ 0 \end{pmatrix}}_{h_2},
  \qquad s,t\in\K.
\]
Here $\rank(A) = 3$, $n = 5$, and $\dim\Ker(A) = 5 - 3 = 2$.
\end{example}


% ============================================================
\section{Geometric interpretation}\label{sec:geometric-interp}
% ============================================================

\index{geometric interpretation!of linear systems}

\subsection{Two equations in two unknowns}

A single linear equation $a_1 x_1 + a_2 x_2 = b$ in two unknowns
describes a \emph{line} in~$\R^2$. A system of two such equations
corresponds to two lines, and the solution set is their intersection.

\begin{center}
\begin{tikzpicture}[>=Stealth, scale=1.0]
  % Unique solution
  \begin{scope}[shift={(0,0)}]
    \draw[->] (-0.5,0) -- (3.5,0) node[right] {$x_1$};
    \draw[->] (0,-0.5) -- (0,3) node[above] {$x_2$};
    \draw[thick, blue!70!black, domain=-0.3:3.2] plot (\x, {2.5 - 0.7*\x});
    \draw[thick, red!70!black, domain=-0.3:3.2] plot (\x, {0.3 + 0.6*\x});
    \fill[black] (1.692, 1.315) circle (2pt);
    \node[below] at (1.75, -0.8) {\small Unique solution};
    \node[below] at (1.75, -1.2) {\small $\rank(A) = 2$};
  \end{scope}
  % No solution
  \begin{scope}[shift={(5.5,0)}]
    \draw[->] (-0.5,0) -- (3.5,0) node[right] {$x_1$};
    \draw[->] (0,-0.5) -- (0,3) node[above] {$x_2$};
    \draw[thick, blue!70!black, domain=-0.3:3.2] plot (\x, {2.5 - 0.7*\x});
    \draw[thick, red!70!black, domain=-0.3:3.2] plot (\x, {1.2 - 0.7*\x});
    \node[below] at (1.75, -0.8) {\small No solution};
    \node[below] at (1.75, -1.2) {\small (parallel lines)};
  \end{scope}
  % Infinitely many
  \begin{scope}[shift={(11,0)}]
    \draw[->] (-0.5,0) -- (3.5,0) node[right] {$x_1$};
    \draw[->] (0,-0.5) -- (0,3) node[above] {$x_2$};
    \draw[thick, blue!70!black, domain=-0.3:3.2] plot (\x, {2.5 - 0.7*\x});
    \draw[thick, red!70!black, dashed, domain=-0.3:3.2]
      plot (\x, {2.5 - 0.7*\x});
    \node[below] at (1.75, -0.8) {\small Infinitely many};
    \node[below] at (1.75, -1.2) {\small (coincident lines)};
  \end{scope}
\end{tikzpicture}
\captionof{figure}{Three possibilities for a $2\times 2$ linear
  system over $\R$.}
\label{fig:2x2-systems}
\end{center}

\subsection{Three equations in three unknowns}

Each equation $a_1 x_1 + a_2 x_2 + a_3 x_3 = b$ describes a
\emph{plane}\index{plane} in~$\R^3$. A system of three such equations
asks for the common intersection of three planes. The possibilities
are richer than in $\R^2$.

\begin{center}
\begin{tikzpicture}[>=Stealth, scale=0.95]
  % --- Unique point ---
  \begin{scope}[shift={(0,0)}]
    % Three planes meeting at a point (schematic)
    \fill[blue!15, opacity=0.7] (-1.2,-0.8) -- (1.5,-0.5)
      -- (1.8,1.3) -- (-0.9,1.0) -- cycle;
    \fill[red!15, opacity=0.6] (-0.5,-1.0) -- (1.8,-0.3)
      -- (1.2,1.5) -- (-1.1,0.8) -- cycle;
    \fill[green!15, opacity=0.5] (-1.0,-0.4) -- (1.0,-1.0)
      -- (1.6,0.9) -- (-0.4,1.5) -- cycle;
    \fill[black] (0.5,0.3) circle (2.5pt);
    \node[below] at (0.4,-1.3) {\small Unique point};
  \end{scope}
  % --- Line ---
  \begin{scope}[shift={(4.5,0)}]
    \fill[blue!15, opacity=0.7] (-1.2,-1.0) -- (1.5,-0.8)
      -- (1.8,1.0) -- (-0.9,0.8) -- cycle;
    \fill[red!15, opacity=0.6] (-1.0,-0.8) -- (1.8,-0.3)
      -- (1.3,1.2) -- (-1.5,0.7) -- cycle;
    \fill[green!15, opacity=0.5] (-0.8,-1.0) -- (1.6,-0.5)
      -- (1.0,1.3) -- (-1.4,0.8) -- cycle;
    \draw[thick, black] (-0.6,-0.6) -- (1.2,0.9);
    \node[below] at (0.4,-1.3) {\small Line};
  \end{scope}
  % --- Plane ---
  \begin{scope}[shift={(9,0)}]
    \fill[blue!20, opacity=0.8] (-1.4,-1.0) -- (1.8,-0.8)
      -- (1.6,1.2) -- (-1.6,1.0) -- cycle;
    \node at (0.1,0.1) {\small (3 coincident)};
    \node[below] at (0.2,-1.3) {\small Plane};
  \end{scope}
  % --- Empty ---
  \begin{scope}[shift={(13.5,0)}]
    \fill[blue!15, opacity=0.6] (-1.2,-0.2) -- (1.5,0.0)
      -- (1.5,0.8) -- (-1.2,0.6) -- cycle;
    \fill[red!15, opacity=0.6] (-1.2,-0.7) -- (1.5,-0.5)
      -- (1.5,0.3) -- (-1.2,0.1) -- cycle;
    \fill[green!15, opacity=0.5] (-1.2,0.4) -- (1.5,0.6)
      -- (1.5,1.4) -- (-1.2,1.2) -- cycle;
    \node[below] at (0.2,-1.3) {\small No solution};
    \node[below] at (0.2,-1.7) {\small (parallel planes)};
  \end{scope}
\end{tikzpicture}
\captionof{figure}{Selected configurations for three planes
  in~$\R^3$: unique intersection point, intersection along a line,
  three coincident planes, and three parallel planes.}
\label{fig:3x3-systems}
\end{center}


% ============================================================
\section{Applications}\label{sec:applications-systems}
% ============================================================

\subsection{Network flows}\label{subsec:network-flows}
\index{network flow}

Consider a network of one-way streets or pipelines with known flow
rates entering and leaving the network. At each junction
(node)\index{node}, conservation of flow requires that the total
flow in equals the total flow out. This yields one linear equation per
node.

\begin{example}{Traffic flow}{ex:traffic-flow}
Consider the network with four intersections where external flows are
given and the internal flows $x_1, x_2, x_3, x_4$ are unknown:

\begin{center}
\begin{tikzpicture}[>=Stealth, node distance=2.5cm,
  every node/.style={circle, draw, minimum size=8mm, inner sep=1pt}]
  \node (A) at (0,2) {$A$};
  \node (B) at (3,2) {$B$};
  \node (C) at (3,0) {$C$};
  \node (D) at (0,0) {$D$};
  % Internal edges
  \draw[->, thick] (A) -- node[above] {$x_1$} (B);
  \draw[->, thick] (B) -- node[right] {$x_2$} (C);
  \draw[->, thick] (C) -- node[below] {$x_3$} (D);
  \draw[->, thick] (D) -- node[left]  {$x_4$} (A);
  % External flows
  \draw[->, thick, blue!70!black] (-1.5,2) -- (A)
    node[midway, above, draw=none, fill=none] {$100$};
  \draw[->, thick, blue!70!black] (B) -- (4.5,2)
    node[midway, above, draw=none, fill=none] {$80$};
  \draw[->, thick, blue!70!black] (4.5,0) -- (C)
    node[midway, above, draw=none, fill=none] {$50$};
  \draw[->, thick, blue!70!black] (D) -- (-1.5,0)
    node[midway, above, draw=none, fill=none] {$70$};
\end{tikzpicture}
\end{center}

Flow conservation at each node:
\[
  \begin{aligned}
    A&: & 100 + x_4 &= x_1, \\
    B&: & x_1 &= 80 + x_2, \\
    C&: & x_2 + 50 &= x_3, \\
    D&: & x_3 &= 70 + x_4.
  \end{aligned}
  \qquad\Longleftrightarrow\qquad
  \begin{aligned}
    x_1 - x_4 &= 100,\\
    x_1 - x_2 &= 80,\\
    x_2 - x_3 &= -50,\\
    x_3 - x_4 &= 70.
  \end{aligned}
\]
Row reduction shows $\rank(A) = 3$ and $n=4$, giving one free
variable. Setting $x_4 = t$:
\[
  x_1 = 100+t, \quad x_2 = 20+t, \quad x_3 = 70+t, \quad x_4 = t.
\]
Physical constraints ($x_i \geq 0$) require $t \geq 0$.
\end{example}


\subsection{Electrical circuits}\label{subsec:circuits}
\index{Kirchhoff's laws}\index{electrical circuit}

In an electrical circuit, Kirchhoff's current law (KCL) states that
the algebraic sum of currents at any node is zero, and Kirchhoff's
voltage law (KVL) states that the algebraic sum of voltage drops
around any closed loop is zero. Combined with Ohm's law
($V = IR$)\index{Ohm's law}, these yield a linear system for the
unknown currents.

\begin{example}{A simple resistive circuit}{ex:circuit}
Consider a circuit with two loops sharing a branch. Let $I_1$, $I_2$,
$I_3$ be the branch currents and $R_1 = 2\,\Omega$,
$R_2 = 4\,\Omega$, $R_3 = 6\,\Omega$ the resistances, with a
voltage source $V = 12$\,V. Applying KCL and KVL:
\[
  \left\{
  \begin{aligned}
    I_1 - I_2 - I_3 &= 0 && (\text{KCL at node}),\\
    2I_1 + 6I_3 &= 12 && (\text{KVL, loop 1}),\\
    4I_2 - 6I_3 &= 0 && (\text{KVL, loop 2}).
  \end{aligned}
  \right.
\]
Solving by Gaussian elimination gives
$I_1 = \frac{36}{11}$\,A,\;
$I_2 = \frac{18}{11}$\,A,\;
$I_3 = \frac{18}{11} \cdot \frac{2}{3} = \frac{12}{11}$\,A.
\end{example}


\subsection{Preview: least squares}\label{subsec:least-squares-preview}
\index{least squares!preview}

When a system $Ax = b$ is \emph{overdetermined} (more equations than
unknowns, typically inconsistent), one seeks the vector $\hat{x}$
that minimises $\norm{Ax - b}^2$. A fundamental result (to be proved
in the chapter on inner product spaces) states that $\hat{x}$
satisfies the \emph{normal equations}\index{normal equations}
\begin{equation}\label{eq:normal-equations}
  \transpose{A}A\,\hat{x} = \transpose{A}b.
\end{equation}
This is a (consistent) square linear system, solvable by Gaussian
elimination. Least squares is the mathematical foundation of linear
regression, signal processing, and data fitting.


% ============================================================
\section{Exercises}\label{sec:exercises-systems}
% ============================================================

% --- Exercise 1 ---
\begin{exercise}{Row reduction practice \basic}{exo:row-reduce-1}
Reduce the following matrix to RREF and find $\rank(A)$:
\[
  A = \begin{pmatrix}
    1 & 2 & -1 & 0\\
    2 & 4 & 1 & 3\\
    3 & 6 & 2 & 5
  \end{pmatrix}.
\]
\end{exercise}

% --- Exercise 2 ---
\begin{exercise}{Solving a $3\times 3$ system \basic}{exo:solve-3x3}
Solve the system
\[
  \left\{
  \begin{aligned}
    x + y + z &= 6,\\
    2x + 3y + z &= 14,\\
    x + 2y - z &= 2.
  \end{aligned}
  \right.
\]
\end{exercise}

% --- Exercise 3 ---
\begin{exercise}{Homogeneous system \basic}{exo:homog-system}
Find the general solution of $Ax = \vzero$ where
\[
  A = \begin{pmatrix}
    1 & 2 & 3\\
    2 & 4 & 6\\
    1 & 2 & 3
  \end{pmatrix}.
\]
What is $\dim\Ker(A)$?
\end{exercise}

% --- Exercise 4 ---
\begin{exercise}{Consistency conditions \basic}{exo:consistency}
For which values of $a\in\R$ is the system
\[
  \left\{
  \begin{aligned}
    x + y &= 1,\\
    2x + 2y &= a,
  \end{aligned}
  \right.
\]
consistent? When consistent, describe the solution set geometrically.
\end{exercise}

% --- Exercise 5 ---
\begin{exercise}{Parametric solutions \intermediate}{exo:parametric-sol}
Solve the system and express the solution in parametric vector form:
\[
  \left\{
  \begin{aligned}
    x_1 + 3x_2 - 2x_3 + x_4 &= 4,\\
    2x_1 + 6x_2 - 3x_3 + 5x_4 &= 11,\\
    x_1 + 3x_2 + x_4 &= 5.
  \end{aligned}
  \right.
\]
Identify the particular solution and a basis for the kernel.
\end{exercise}

% --- Exercise 6 ---
\begin{exercise}{Rank and the Rouch\'{e}--Capelli theorem
\intermediate}{exo:rouche-capelli-app}
Let
\[
  A = \begin{pmatrix}
    1 & 2 & 1\\
    2 & 4 & 3\\
    3 & 6 & 4
  \end{pmatrix},
  \qquad
  b = \begin{pmatrix} 1\\3\\4 \end{pmatrix}.
\]
Compute $\rank(A)$ and $\rank(A\mid b)$. Is $Ax = b$ consistent?
If so, find the solution set.
\end{exercise}

% --- Exercise 7 ---
\begin{exercise}{Family of systems \intermediate}{exo:family-systems}
For which values of $\lambda\in\R$ does the system
\[
  \left\{
  \begin{aligned}
    x + y + \lambda z &= 1,\\
    x + \lambda y + z &= 1,\\
    \lambda x + y + z &= 1,
  \end{aligned}
  \right.
\]
have (a) a unique solution, (b) infinitely many solutions,
(c) no solution?
\end{exercise}

% --- Exercise 8 ---
\begin{exercise}{Network flow \intermediate}{exo:network-flow}
In a road network with five nodes, the flow conservation equations
yield the system
\[
  \left\{
  \begin{aligned}
    x_1 + x_2 &= 200,\\
    x_2 + x_3 &= 150,\\
    x_3 + x_4 &= 180,\\
    x_4 + x_5 &= 120.
  \end{aligned}
  \right.
\]
Find the general solution. If all flows must be non-negative,
determine the feasible range for~$x_1$.
\end{exercise}

% --- Exercise 9 ---
\begin{exercise}{Computing a kernel basis \intermediate}{exo:kernel-basis}
Find a basis for $\Ker(A)$ where
\[
  A = \begin{pmatrix}
    1 & 0 & -2 & 1 & 3\\
    0 & 1 & 3 & -1 & 0\\
    2 & 1 & -1 & 1 & 6\\
    1 & 1 & 1 & 0 & 3
  \end{pmatrix}.
\]
Verify that $\rank(A) + \dim\Ker(A) = 5$.
\end{exercise}

% --- Exercise 10 ---
\begin{exercise}{Rank inequalities \challenging}{exo:rank-ineq}
Let $A\in\Mat_{m,n}(\K)$ and $B\in\Mat_{n,p}(\K)$. Prove:
\begin{enumerate}[label=(\alph*)]
  \item $\rank(AB) \leq \min(\rank(A), \rank(B))$.
  \item \emph{Sylvester's rank inequality:}\index{Sylvester's rank inequality}
    $\rank(AB) \geq \rank(A) + \rank(B) - n$.
\end{enumerate}
\textit{Hint for (b):} use the rank--nullity theorem and show
$\Ker(B) \subseteq \Ker(AB)$.
\end{exercise}

% --- Exercise 11 ---
\begin{exercise}{Normal equations \challenging}{exo:normal-eq}
Consider the inconsistent system (three equations, two unknowns):
\[
  \begin{pmatrix} 1 & 1\\ 1 & 2\\ 1 & 3 \end{pmatrix}
  \begin{pmatrix} x_1\\ x_2 \end{pmatrix}
  = \begin{pmatrix} 1\\ 2\\ 4 \end{pmatrix}.
\]
\begin{enumerate}[label=(\alph*)]
  \item Form the normal equations $\transpose{A}A\,\hat{x} =
    \transpose{A}b$.
  \item Solve for $\hat{x}$ by Gaussian elimination.
  \item Interpret $\hat{x}$ as the coefficients of the best-fit
    line $y = x_1 + x_2 t$ through the data points
    $(1,1)$, $(2,2)$, $(3,4)$.
\end{enumerate}
\end{exercise}

% --- Exercise 12 ---
\begin{exercise}{Intersection of subspaces via systems
\challenging}{exo:intersect-subspaces}
Let $U = \Span\!\left(\begin{pmatrix}1\\1\\0\\1\end{pmatrix},
\begin{pmatrix}0\\1\\1\\0\end{pmatrix}\right)$ and
$W = \Span\!\left(\begin{pmatrix}1\\0\\1\\1\end{pmatrix},
\begin{pmatrix}1\\1\\1\\0\end{pmatrix}\right)$ be subspaces
of~$\R^4$.
\begin{enumerate}[label=(\alph*)]
  \item Set up a linear system whose solutions describe $U\cap W$.
  \item Solve the system and find a basis for $U\cap W$.
  \item Verify the Grassmann formula:
    $\dim(U+W) = \dim U + \dim W - \dim(U\cap W)$.
\end{enumerate}
\end{exercise}

% --- Exercise 13 ---
\begin{exercise}{Proving row equivalence preserves solution sets
\intermediate}{exo:ero-proof}
Give a detailed proof that each elementary row operation on the
augmented matrix $(A\mid b)$ preserves the solution set of $Ax = b$.
Explicitly construct the inverse operation in each case.
\end{exercise}

% --- Exercise 14 ---
\begin{exercise}{RREF uniqueness \challenging}{exo:rref-unique-detail}
Provide a complete proof that two row-equivalent matrices in reduced
row echelon form must be identical.
\textit{Hint:} proceed by induction on the number of columns, using
the fact that the two matrices have identical null spaces.
\end{exercise}


% ============================================================
\section{Chapter summary}\label{sec:summary-ch3}
% ============================================================

\begin{itemize}
  \item A \textbf{system of linear equations}\index{system of linear equations}
    $Ax = b$ involves a coefficient matrix $A\in\Mat_{m,n}(\K)$, an
    unknown vector $x\in\K^n$, and a right-hand side $b\in\K^m$.

  \item \textbf{Elementary row operations}\index{elementary row operations}
    (row swap, row scaling, row replacement) transform a system into a
    row-equivalent one with the same solution set.

  \item \textbf{Gaussian elimination}\index{Gaussian elimination}
    reduces $A$ to \textbf{row echelon form}\index{row echelon form}
    (REF); \textbf{Gauss--Jordan elimination}%
    \index{Gauss--Jordan elimination} continues to the unique
    \textbf{reduced row echelon form}\index{reduced row echelon form}
    (RREF).

  \item The \textbf{rank}\index{rank} $\rank(A)$ is the number of
    pivots, equal to the dimension of both the column space and the
    row space.

  \item \textbf{Rouch\'{e}--Capelli theorem}%
    \index{Rouche-Capelli theorem@Rouch\'{e}--Capelli theorem}:
    $Ax = b$ is consistent $\iff$ $\rank(A) = \rank(A\mid b)$.

  \item The \textbf{general solution}\index{general solution} of a
    consistent system is
    $\cS = x_p + \Ker(A)$,
    an affine subspace of dimension $n - \rank(A)$.

  \item \textbf{Rank--nullity:}\index{rank--nullity theorem}
    $\rank(A) + \dim\Ker(A) = n$ (number of columns).

  \item Geometrically, a linear equation in $\R^2$ defines a line, in
    $\R^3$ a plane; the solution set of a system is the intersection
    of these geometric objects.

  \item Applications include network flows, electrical circuits
    (Kirchhoff's laws), and least squares (normal equations).
\end{itemize}
